\documentclass[a4paper,11pt]{article}
\usepackage{fancyhdr,graphics,graphicx,xy}
\usepackage{blkarray}
\usepackage{ifthen}
\usepackage{amssymb,amsmath,mathtools,color}
\pagestyle{fancy}
\usepackage{xparse}
\usepackage{float}
\usepackage{nicematrix}


%\textwidth = 6.5 in
%\textheight = 9 in
%\oddsidemargin = 0.0 in
%\evensidemargin = 0.0 in
%\topmargin = 0.0 in
%\headheight = 0.0 in
%\headsep = 0.0 in
%\parskip = 0.2in
%\parindent = 0.0in


%%%%%%%%%%%%%%%%%%%%%%%%%%%%%%%%%%%%%%%%%%
%%% General math macros
%%%%%%%%%%%%%%%%%%%%%%%%%%%%%%%%%%%%%%%%%%

\newcommand{\tf}{\tfrac} 
\renewcommand{\Re}{\mathrm{Re}}

\newcommand{\Def}{\coloneqq}
\newcommand{\lgg}{\operatorname{lg}}
\newcommand{\R}{\mathbb{R}}
\newcommand{\C}{\mathbb{C}}
\newcommand{\D}{\mathbb{D}}
\newcommand{\A}{\mathcal{A}}
\newcommand{\F}{\mathcal{F}}
\newcommand{\Q}{\mathbb{Q}}
\newcommand{\N}{\mathbb{N}} 
\newcommand{\Quat}{\mathbb{J}}
\newcommand{\Z}{\mathbb{Z}}
\newcommand{\T}{\mathbb{T}}
\newcommand{\U}{\mathcal{U}}
\newcommand{\LP}{\operatorname{L}}
\newcommand{\abs}[1]{\left| {#1} \right|}
\newcommand{\sgn}{\operatorname{sgn}}

\newcommand{\tendsto}{\to}
\newcommand{\tendto}{\to}

%%%%%%%%%%%%%%%%%%%%%%%%%%%%%%%%%%%%%%%%%%
%%% Probability Macros
%%%%%%%%%%%%%%%%%%%%%%%%%%%%%%%%%%%%%%%%%%
\newcommand{\stdcg}{N_\bbc(0,1)}

\newcommand{\expec}{\mathbb{E}}
\newcommand{\expect}{\mathbb{E}}
\newcommand{\Exp}{\mathbb{E}}
\newcommand{\E}{\mathbb{E}}
\newcommand{\prob}{\mathbb{P}}
\newcommand{\filt}{\mathscr{F}}
\newcommand{\Prob}{\prob}
\renewcommand{\Pr}{\prob}
\newcommand{\Ent}{\operatorname{Ent}}
\newcommand{\Var}{\operatorname{Var}}
\newcommand{\Cov}{\operatorname{Cov}}
\newcommand{\dtv}{d_{TV}}
\newcommand{\weakto}{\Rightarrow}
\newcommand{\lawof}{\mathscr{L}}
\newcommand{\Bernoulli}{\operatorname{Bernoulli}}
\newcommand{\Binomial}{\operatorname{Binom}}
\newcommand{\Poisson}{\operatorname{Poisson}}
\newcommand{\Normal}{\operatorname{Normal}}
\newcommand{\Unif}{\operatorname{Unif}}
\newcommand{\As}{\ensuremath{\text{a.s.}}}
\newcommand{\distequals}{\overset{\mathrm{d}}{=}}


\newcommand{\asstnumber}{5}
\newcommand{\assigntop}{
	\setlength{\headheight}{33.40004pt}
  \fancyhead[L]{}
  \fancyhead[C]{\Large Math 356 -- Fall 2023\\ Assignment \asstnumber}
  \fancyhead[R]{ }
}

\newtheorem{theorem}{Theorem}
\newtheorem{corollary}[theorem]{Corollary}
\newtheorem{definition}{Definition}
\newcommand{\p}[1]{\mathbb{P}(#1)}
\newcommand{\px}[1]{\mathbb{P}_x(#1)}
\newcommand{\py}[1]{\mathbb{P}_y(#1)}
\newcommand{\I}[1]{\mathbf{1}_{#1}}
\newcommand{\e}{\mathbb{E}}

\begin{document}
\assigntop

\section*{Question 1}

\subsection*{(a)} 
Let $G$ be a $3$-regular graph. Show that if $G$ contains a Hamiltonian cycle then $\chi' (G) = 3$.
Let $C$ be a hamiltonian cycle in $G$, that means $C$ is a path with the same vertex endpoints and $V(C) = V(G)$. Moreover, $C$ is even, since, by the handshaking lemma 
$$ \sum_{v \in V(G)} \deg (v) = 3|V(G)| = 2 \cdot E(G) \implies |V(G)| \text{ is even } $$
Then there exists a coloring with $E(C) \to \{1,2\}$, and since no edges of $G-C$ can be incident to the same vertex, by $3$-regularity of $G$, then we can color all edges in $G - C$ with color $3$, which yields a $3$-edge coloring of $G$.
\subsection*{(b)} 
From the lectures on Hamiltonian cycles, we have the canonical counter example of a graph with no Hamiltonian cycle from lecture 7.
\begin{figure}[H]
	\centering
	\includegraphics[width=7cm]{ex1.jpeg}
	\caption{Counter example}
	\label{fig:ex1}
\end{figure}
By observation, the graph is $3$-regular.

\section*{(b)}
We show first that graphs with no $K_4$ minor are $k$-degenerate. Let $H$ be a subgraph of $G$, then if $H = G$, by lemma 13.2, we have that $\exists v \in V(G)$ such that $\deg(v) \leq 2$. If $H \neq G$, then select $X_1 \subseteq G-H$ to be a clique, where $|X_1| \leq 2$ in $G - H$, then there exists $v \in G-X$ such that $\deg (v) \leq 2$. If $v \in H$ then we are done, otherwise, consider $G - X_1$, which doesn't have a $K_4$ minor, since $G$ doesn't have a $K_4$ minor. Select $X_i$ such that $X_i \subseteq G - X_1 - X_2 \ldots -X_{i-1} - H$, where $G - X_1 - X_2 - \ldots -X_{i-1}$ doesn't have a $K_4$ minor, until we get a $v \in H$ such that $\deg (v) \leq 2$. 
\end{document}
